%!TEX root = ../dissertation.tex
\chapter{Theoretical Background}\label{chapt:background}

% =====================================================================
%  DRAFT v2 (Asma) -- full rewrite. ~13 pages target.
%  This chapter is TEXTBOOK material: define the tools, no results,
%  no competing CIR systems (those live in Chapter~\ref{chapt:related}).
%  HIGHLIGHT markers = optional figures or things to double-check.
% =====================================================================

This chapter develops the technical vocabulary on which the rest of the thesis
rests. The goal is deliberately narrow: to \emph{define the tools}, not to
survey competing systems (which Chapter~\ref{chapt:related} does) or to report
any results. We build the material bottom-up. Section~\ref{sect:bg_retrieval}
fixes the fundamentals of content-based image retrieval -- embeddings,
similarity, nearest-neighbour search, and how rankings are scored.
Section~\ref{sect:bg_transformer} introduces the Transformer and the Vision
Transformer, the architecture shared by every encoder and language model used
later. Section~\ref{sect:bg_clip} describes contrastive vision--language
pretraining through CLIP, and Section~\ref{sect:bg_siglip} its SigLIP and
SigLIP2 successors. Section~\ref{sect:bg_mllm} introduces multimodal large
language models, the generative component of our pipeline.
Section~\ref{sect:bg_cir} then frames composed image retrieval as a task and
contrasts early- with late-fusion designs, and
Section~\ref{sect:bg_postproc} closes with the training-free feature
post-processing used by the \textsc{Basic} baseline.

% =====================================================================
\section{Image Retrieval Fundamentals}\label{sect:bg_retrieval}
% =====================================================================

\subsection{Embeddings and Similarity}

Modern content-based image retrieval represents every image as a fixed-length
vector, or \emph{embedding}, produced by a learned encoder
$\enc : \mathcal{I} \to \mathbb{R}^d$ that maps the raw image space
$\mathcal{I}$ into a $d$-dimensional space in which semantically similar inputs
lie close together. Retrieval then reduces to nearest-neighbour search: a query
is encoded into the same space, and the database is ranked by similarity to the
query embedding. The quality of a retrieval system is therefore largely
inherited from the geometry of the embedding space its encoder induces, which
is why the choice of encoder (Sections~\ref{sect:bg_clip}--\ref{sect:bg_siglip})
is the dominant design decision in this thesis.

The standard similarity measure is the cosine similarity, which compares the
\emph{directions} of two vectors irrespective of their magnitude,
\begin{equation}\label{eq:cosine}
  \simop(u, v) = \frac{u^{\top} v}{\lVert u \rVert\, \lVert v \rVert}.
\end{equation}
When all embeddings are $\ell_2$-normalised to unit length, as is standard for
the encoders used here, the cosine similarity equals the dot product
$u^{\top}v$, and ranking by cosine similarity is equivalent to ranking by
(negative) squared Euclidean distance, since
$\lVert u - v\rVert^2 = 2 - 2\,u^{\top}v$ for unit vectors. Throughout the
thesis, ``similarity'' means the cosine similarity between $\ell_2$-normalised
embeddings unless stated otherwise.

\subsection{Nearest-Neighbour Search}

Given a query embedding $q$ and a database
$\DB = \{\xdb_1, \dots, \xdb_N\}$ of $N$ image embeddings, retrieval returns the
database images sorted by decreasing similarity $\simop(q, \xdb_i)$. For the
database sizes considered in this thesis ($N$ up to ${\sim}752$K), an exact
search expressed as a single matrix--vector product is both feasible and exact,
so we use brute-force search and do not rely on approximate nearest-neighbour
indexing; the practical cost of this choice is quantified in
Chapter~\ref{chapt:results}. The output of a query is a full ranking of the
database, from which all evaluation metrics are computed.

\subsection{Evaluation: Precision, Recall, and Average Precision}

Retrieval quality is measured against a set of \emph{relevant} (positive)
database items per query. The two elementary quantities are \emph{precision}
(the fraction of retrieved items that are relevant) and \emph{recall} (the
fraction of relevant items that are retrieved). Evaluated at a cutoff $k$,
$\mathrm{Recall}@k$ is the fraction of the relevant items found within the
top-$k$ results, and $\mathrm{Hit}@k$ is $1$ if at least one relevant item
appears in the top-$k$ and $0$ otherwise. These cutoff metrics are intuitive
but discard the order \emph{within} the cutoff.

To summarise an entire ranking with a single order-sensitive number, retrieval
uses \emph{average precision} (AP), the area under the precision--recall curve,
which rewards rankings that place relevant items near the top. For a query with
relevant set $R$,
\begin{equation}\label{eq:ap_basic}
  \mathrm{AP} = \frac{1}{|R|} \sum_{k=1}^{N} P(k)\,\mathbb{1}[\,\text{item at
  rank } k \text{ is relevant}\,],
\end{equation}
where $P(k)$ is the precision at cutoff $k$. The \emph{mean average precision}
(mAP) averages AP over queries. The exact trapezoidal variant used by the i-CIR
benchmark, together with the distinction between macro- and micro-averaging, is
defined where it is first needed in Section~\ref{sect:method_eval}.

% =====================================================================
\section{Transformers and Vision Transformers}\label{sect:bg_transformer}
% =====================================================================

Every encoder and language model used in this thesis is built from the same
backbone: the Transformer~\cite{vaswani2017attention}. Understanding its core
mechanism, self-attention, clarifies why a single architecture can serve as an
image encoder, a text encoder, and the body of a large language model alike.

\subsection{Self-Attention and the Transformer}

A Transformer operates on a \emph{sequence} of token vectors. Its defining
operation, \emph{scaled dot-product attention}, lets every token aggregate
information from all other tokens, with weights that depend on content rather
than position. Each token is projected into a \emph{query}, a \emph{key}, and a
\emph{value} vector; stacking these over the sequence gives matrices $Q$, $K$,
and $V$, and the attention output is
\begin{equation}\label{eq:attention}
  \mathrm{Attention}(Q, K, V) =
  \mathrm{softmax}\!\left(\frac{Q K^{\top}}{\sqrt{d_k}}\right) V,
\end{equation}
where $d_k$ is the key dimension and the $\sqrt{d_k}$ factor stabilises the
softmax. Intuitively, the dot product $Q K^{\top}$ measures how relevant each
token is to every other, the softmax turns these scores into a distribution,
and the output is a content-weighted average of the value vectors. In practice
several attention ``heads'' run in parallel on different learned projections
(\emph{multi-head attention}), letting the model attend to different kinds of
relations simultaneously.

A Transformer block wraps multi-head self-attention and a position-wise
feed-forward network, each with a residual connection and layer normalisation;
blocks are stacked to form deep models. Because attention is permutation-
invariant, explicit \emph{positional encodings} are added to the token
embeddings so the model can use order. The architecture has two crucial
properties for this thesis: it is highly parallelisable and it scales
gracefully, which is precisely why training on web-scale image--text and text
corpora -- the regime that produces both CLIP-style encoders and large language
models -- became practical.

\subsection{The Vision Transformer}

The Vision Transformer (ViT)~\cite{dosovitskiy2021vit} applies this machinery to
images with minimal change. An image is split into a grid of fixed-size,
non-overlapping \emph{patches} (for example $16{\times}16$ pixels); each patch
is flattened and linearly projected into a token embedding, and a learned
positional encoding is added. The resulting sequence of patch tokens -- usually
together with a special learnable \texttt{[CLS]} token -- is processed by a
stack of Transformer blocks, after which a single image-level embedding is read
out, either from the \texttt{[CLS]} token or by pooling the patch tokens. Self-
attention thus lets every patch interact with every other, giving the model a
global receptive field already in its first layer.

The ViT is the visual encoder of the contrastive models studied here and, in
the form of a vision encoder feeding a language model, the visual front-end of
the multimodal LLMs of Section~\ref{sect:bg_mllm}. The remainder of the chapter
treats these Transformer encoders as black boxes that map an image or a text to
an embedding; what differs between them is not the architecture but the
objective they were trained with, to which we now turn.

% =====================================================================
\section{Contrastive Vision--Language Pretraining: CLIP}\label{sect:bg_clip}
% =====================================================================

The encoders used in this thesis come from the family of \emph{contrastive
vision--language models}, of which CLIP (Contrastive Language--Image
Pretraining)~\cite{radford2021clip} is the canonical example. CLIP is a
\emph{dual-encoder} -- also called \emph{dual-tower} or \emph{two-tower} --
model: a vision encoder $\encv$ (a ViT or a ResNet) maps an image to an
embedding, and a \emph{separate} text encoder $\enct$ (a Transformer) maps a
text string to an embedding in the \emph{same} $d$-dimensional space, with both
embeddings $\ell_2$-normalised. The two towers share no weights and never fuse
internally; they interact only through the geometry of the shared output space,
where matching is a simple dot product. This separation is what makes the
architecture so well suited to retrieval: the (large) image side can be encoded
once, offline, and indexed, while only the query is encoded at search time.

\subsection{The Contrastive Objective}

CLIP is trained on a very large collection of image--text pairs harvested from
the web (400M pairs for the original model; later open reproductions use the
LAION datasets~\cite{schuhmann2022laion5b,ilharco2021openclip}). Training uses a
symmetric \emph{contrastive} objective of the InfoNCE form: within a batch of
$B$ image--text pairs, the model treats the $B$ true pairings as positives and
all $B^2 - B$ mismatched pairings as negatives, maximising the similarity of
the former relative to the latter. With image embeddings $u_i$, text embeddings
$v_j$, and a learned temperature $\tau$, the image-to-text loss is
\begin{equation}\label{eq:infonce}
  \mathcal{L}_{\text{i2t}} = -\frac{1}{B}\sum_{i=1}^{B}
  \log \frac{\exp(u_i^{\top} v_i / \tau)}
            {\sum_{j=1}^{B}\exp(u_i^{\top} v_j / \tau)},
\end{equation}
and the total loss symmetrises this with the analogous text-to-image term. The
softmax in Equation~\eqref{eq:infonce} is normalised over the whole batch, so
the quality of the learned space improves with batch size, which has practical
consequences revisited in Section~\ref{sect:bg_siglip}.

\subsection{The Shared Space and Zero-Shot Transfer}

The property this thesis depends on is the resulting \emph{shared embedding
space}: an image and a text that describes it are mapped to nearby vectors.
This directly enables zero-shot cross-modal retrieval -- a text query can be
matched against image embeddings without any task-specific training -- and it
is the substrate on which every retrieval signal in this thesis (image, text,
and generated caption) is compared.

The same alignment underlies CLIP's celebrated \emph{zero-shot transfer}.
\emph{Zero-shot} learning refers to performing a task for which the model
received no labelled training examples; CLIP achieves it by recasting the task
in language. To classify an image, for instance, one embeds candidate class
names as text prompts (``a photo of a \{class\}'') and returns the class whose
text embedding is most similar to the image embedding -- no classifier is
trained. Because the alignment was learned from open-domain web data, this
transfers to new datasets and label sets out of the box. This capability is the
conceptual foundation of the zero-shot CIR setting in which this thesis
operates: it lets pretrained, frozen encoders be repurposed for composed
retrieval without any CIR-specific supervision.

% HIGHLIGHT (optional): a CLIP dual-encoder diagram would help here.
% \begin{figure}[t] \centering
%   \includegraphics[width=0.85\textwidth]{figures/clip.png}
%   \caption{The CLIP dual-encoder architecture.}\label{fig:clip} \end{figure}

% =====================================================================
\section{SigLIP and SigLIP2}\label{sect:bg_siglip}
% =====================================================================

SigLIP (Sigmoid Loss for Language--Image Pretraining)~\cite{zhai2023siglip}
keeps the CLIP dual-encoder architecture but replaces the softmax-based
contrastive loss of Equation~\eqref{eq:infonce} with a pairwise \emph{sigmoid}
loss. Each image--text pair is treated as an independent binary classification
problem -- matching or not -- so the loss does not require a softmax normalised
over the whole batch:
\begin{equation}\label{eq:siglip}
  \mathcal{L} = -\frac{1}{B}\sum_{i=1}^{B}\sum_{j=1}^{B}
  \log \sigma\!\big(y_{ij}\,(t\, u_i^{\top} v_j + b)\big),
\end{equation}
where $y_{ij} = +1$ for matching pairs and $-1$ otherwise, $\sigma$ is the
logistic sigmoid, and $t, b$ are a learned scale and bias. Decoupling the loss
from batch-wide normalisation makes training more stable, more memory-efficient,
and effective at smaller batch sizes, which in practice yields encoders with
stronger image--text alignment than their CLIP counterparts trained under
comparable budgets.

SigLIP2~\cite{tschannen2025siglip2} extends SigLIP with additional training
objectives -- including self-distillation and captioning-style losses -- and
multilingual data, improving semantic understanding, localisation, and dense
features. For this thesis the relevant consequence is empirical: SigLIP2
provides a particularly well-aligned shared embedding space, which (as
Chapter~\ref{chapt:results} shows) makes it the strongest backbone for our
pipeline. The thesis treats CLIP, SigLIP, and SigLIP2 uniformly as frozen
dual encoders that supply $\encv$ and $\enct$; their differences are confined to
\emph{how} that shared space was trained, not to how it is used downstream.

% =====================================================================
\section{Multimodal Large Language Models}\label{sect:bg_mllm}
% =====================================================================

The second model family this thesis relies on is generative rather than
contrastive. A \emph{large language model} (LLM) is a Transformer trained
autoregressively on massive text corpora to predict the next token; at scale,
this single objective yields models that follow instructions and produce fluent,
context-dependent text~\cite{brown2020gpt3,vaswani2017attention}. A
\emph{multimodal large language model} (MLLM) -- also called a vision--language
model in this generative sense -- augments an LLM with a visual input pathway,
so that it can take an image (or several images) together with a text prompt and
generate free-form text conditioned on both.

A typical MLLM couples a pretrained vision encoder (a ViT,
Section~\ref{sect:bg_transformer}) to a pretrained LLM through a lightweight
\emph{connector} -- for example an MLP projector or a set of learned query
tokens -- that maps visual features into the LLM's token-embedding space; the
combined model is then instruction-tuned on multimodal data. The InternVL
family~\cite{wang2025internvl35} used in this thesis follows exactly this
recipe, pairing a strong vision encoder with an LLM and supporting
high-resolution images through a tiling scheme that splits an image into several
sub-images. Qwen-based vision--language models~\cite{bai2025qwen3vl} are
architecturally similar and are used here as a comparison point.

The capability this thesis exploits is \emph{instruction-conditioned visual
description}: given the reference image and a prompt that states the desired
modification, an MLLM can be asked to describe the \emph{intended target} image,
i.e.\ to apply the edit in language. This is fundamentally different from a
contrastive dual encoder. A dual encoder can only \emph{embed} a given text or
image; it cannot \emph{compose} a new description from an image and an
instruction. The MLLM, being generative and instruction-following, performs
exactly this composition, producing a self-contained caption that a
dual-encoder text branch can subsequently embed and use for retrieval. Crucially,
the thesis uses the MLLM purely as a frozen, off-the-shelf caption generator --
it is never fine-tuned -- so the overall pipeline remains zero-shot with respect
to the retrieval task.

% =====================================================================
\section{Composed Image Retrieval as a Task}\label{sect:bg_cir}
% =====================================================================

Composed image retrieval (CIR) generalises the two classical retrieval modes.
\emph{Content-based image retrieval} uses an image as the query, and
\emph{text-to-image retrieval} uses a text; CIR uses both at once. The query is
a pair $(\qv, \qt)$ of a reference image and a modification text, and the system
must retrieve the image that results from applying the modification to the
reference~\cite{vo2019tirg}. The defining challenge, introduced in
Chapter~\ref{chapt:intro}, is that $\qt$ expresses a \emph{relative} change --
``the same bag but in red'', ``without the people'' -- and is therefore not, on
its own, a description of the target. A system must read the reference image to
know what is being modified and the text to know how, and combine the two.

How that combination is performed splits the field into two families.
\emph{Early fusion} merges the modalities \emph{before} matching, producing a
single composed query embedding -- typically with a trained fusion network that
takes $(\encv(\qv), \enct(\qt))$ and outputs one vector, which is then matched
against database image embeddings~\cite{vo2019tirg}. \emph{Late fusion} keeps
the modalities separate: it computes independent similarity scores per modality
against each candidate and combines the \emph{scores}, for example by a weighted
sum or a product. Late fusion needs no trained combiner, which makes it the
natural choice in the zero-shot setting; the \textsc{Basic} baseline and the
method proposed in this thesis are both late-fusion methods
(Chapters~\ref{chapt:related} and~\ref{chapt:method}).

A second, orthogonal axis is supervision. \emph{Supervised} CIR trains on
annotated triplets (reference, modification, target) drawn from a CIR training
set~\cite{vo2019tirg,liu2021cirr}, whereas \emph{zero-shot} CIR performs no
CIR-specific training and instead reuses frozen, pretrained
encoders~\cite{saito2023pic2word,baldrati2023circo}. This thesis operates
entirely in the zero-shot, late-fusion regime; the specific systems that
populate these categories are reviewed in Chapter~\ref{chapt:related}.

% =====================================================================
\section{Feature Post-Processing for Retrieval}\label{sect:bg_postproc}
% =====================================================================

Even with a fixed encoder, the geometry of an embedding space can be reshaped
for retrieval by simple, training-free post-processing. The \textsc{Basic}
baseline (Section~\ref{sect:method_basic}) combines several such techniques;
this section defines them in the abstract, deferring their assembly and
empirical effect to Chapters~\ref{chapt:method} and~\ref{chapt:results}.

\paragraph{Centering.} Embeddings produced by contrastive encoders often share
a dominant common component -- an ``embedding-space bias'' -- that inflates the
similarity of \emph{all} pairs and reduces the discriminativeness of the cosine
similarity. \emph{Centering} subtracts a fixed mean vector $\mu$ from every
embedding and re-normalises,
\begin{equation}\label{eq:centering}
  \tilde{\enc}(x) = \frac{\enc(x) - \mu}{\lVert \enc(x) - \mu \rVert},
\end{equation}
which removes this shared component and typically sharpens retrieval. The mean
$\mu$ can be estimated from the database itself (a global mean) or from an
external reference corpus such as a sample of LAION~\cite{schuhmann2022laion5b};
the effect of this choice is studied empirically in
Section~\ref{sect:res_centering}.

\paragraph{Dimensionality reduction and whitening.} Principal component analysis
(PCA) projects embeddings onto their leading directions of variance and, when
combined with \emph{whitening}, equalises the variance across those directions.
In retrieval this can suppress noisy or redundant dimensions and has a long
history as a post-processing step~\cite{radenovic2018finetuning}.

\paragraph{Query expansion.} Query expansion refines a query by aggregating it
with its own top-ranked neighbours: after a first retrieval pass, the query
embedding is combined with the embeddings of the highest-scoring database items
and the search is repeated. This propagates information from confident matches
and is a classic technique in instance
retrieval~\cite{philbin2007oxford,radenovic2018finetuning}.

These operations are entirely unsupervised and add no learned parameters, which
is why they fit naturally into a zero-shot pipeline. Chapter~\ref{chapt:method}
describes how \textsc{Basic} assembles them into a retrieval system, and
Chapter~\ref{chapt:results} quantifies what each contributes.
